% =====================================================================
%  SHERLOCK manuscript — Results sections 2.5–2.6 (Norman 2019,
%  combinatorial perturbations and genetic interactions)
%
%  FIGURE REFERENCE NOTES (update labels to match the master file):
%   * \ref{fig:fig6}  = NORMAN MAIN FIGURE. Panels assumed, in order:
%       (a) SHERLOCK latent space (UMAP of z; singles + combinations,
%           coloured by Norman Supplementary Table 6 gene module)
%       (b) Counterfactual ATE (SHERLOCK / GEARS / cVAE / sVAE /
%           contrastiveVI / scGen; all, singles, combos, held-out combos)
%       (c) Norman GI labels vs SHERLOCK GI prediction (observed mode)
%           — GI-metric embedding
%       (d) GI classification (observed mode, all labelled pairs):
%           per-class AUC-ROC + F1/ARI/AMI
%       (e) GI categories in reconstructed gene space
%     -> If panel letters differ in the master file, only the
%        \ref{fig:fig6}X letters below need editing.
%   * Supplementary figures cited as: fig:SuppNormanCov,
%     fig:SuppNormanLatent, fig:SuppNormanRepro, fig:SuppNormanGears,
%     fig:SuppNormanGIheldout, fig:SuppNormanGIraw. Renumber to the
%     \label{fig:SupplementaryfigN} scheme used elsewhere when merging.
%   * NEW citation keys used below — add to refs.bib if absent:
%     norman2019exploring, roohani2023gears, akiba2019optuna, sen1968estimates,
%     horlbeck2018mapping, cuenda2007p38, friedman2007cebp,
%     rosenbauer2007transcription, davis2003eto, siatecka2011klf1,
%     sherr1999cdk, caunt2013dual, mohapatra2013cbl, tsygankov2013tula,
%     youle2008bcl2.
% =====================================================================

\subsection{SHERLOCK models combinatorial perturbations as compositional latent interventions}

The three preceding sections addressed perturbation grouping, downstream target-gene inference, and
condition-dependent responses. The fourth dimension of a perturbation screen---how perturbations
interact when applied jointly---requires a model in which a combination is not an independent label
but a composition of the mechanisms of its constituents. We therefore evaluated SHERLOCK on the
CRISPRa combinatorial screen of Norman et al. \cite{norman2019exploring}, which profiles 105 single
and 131 double perturbations in K562 cells (111,255 cells; 3,270 genes retained after keeping genes
with mean expression $>0.5$ UMI per cell together with all perturbation-target genes). K562 is a
chronic myelogenous leukaemia line poised between erythroid and megakaryocytic fates, and the
perturbation library was constructed around regulators that push it along those axes---erythroid
(\textit{KLF1}, \textit{LYL1}, \textit{SLC4A1}), granulocytic/monocytic (\textit{CEBPA},
\textit{CEBPB}, \textit{CEBPE}, \textit{SPI1}), growth-factor and MAPK signalling (\textit{MAPK1},
\textit{MAP2K3}, \textit{MAP2K6}, \textit{DUSP9}), and cell-cycle and apoptotic control
(\textit{CDKN1A/B/C}, \textit{BAK1}, \textit{BCL2L11})---so that combinations probe how these
programs compose. Eighty-eight of the double perturbations carry expert-assigned
genetic-interaction (GI) annotations spanning eight categories.

In SHERLOCK a combination label $A{+}B$ is decomposed into its components, and the joint latent
intervention is constructed from the two single-perturbation shifts rather than from a free
embedding of the pair:
\[
\Delta_{AB} \;=\; c_A(\rho_A,\rho_B)\,\Delta_A \;+\; c_B(\rho_A,\rho_B)\,\Delta_B \;+\; s(\rho_A,\rho_B),
\qquad \Delta_p = W_p \odot \big(\Sigma^{1/2}\rho_p\big).
\]
The parent coefficients $c_A,c_B$ are scalar reweightings of each parent's shift, produced by a
shared linear map applied to the ordered pair $[\rho_{\text{self}},\rho_{\text{other}}]$ in both
orderings; this makes the construction swap-equivariant ($c_A(A,B)=c_B(B,A)$, so $A{+}B$ and $B{+}A$
are identical) while leaving the two coefficients free to move independently. Passing the score
through $1+\tanh(\cdot)$ confines each coefficient to $(0,2)$: a parent's program can be damped or up
to doubled, but cannot change sign or diverge. Biologically, this is the degree of freedom that
distinguishes suppression (both parents damped) and potentiation (both amplified) from dominance of
one parent over the other, all of which are attested categories in the Norman annotation. The
interaction residual $s$ is a symmetrised multilayer perceptron over $(\rho_A,\rho_B)$ and is
therefore not restricted to the span of the two parent shifts, so it can express interaction
directions---asymmetric masking, or transcriptional states absent from either single, as in
neomorphic pairs---that a bilinear or low-rank interaction term cannot represent.

Two design choices make additivity the model's null hypothesis rather than an artefact of
parameterisation. First, the parent-coefficient map is initialised near zero so that
$c_A=c_B\approx 1$ at the start of training and $\Delta_{AB}$ begins exactly at $\Delta_A+\Delta_B$;
the model must be driven off additivity by the data. Second, deviations are explicitly priced: the
objective carries an $\ell_2$ penalty on the interaction residual $s$ and on $(c-1)^2$. Because
nothing else in the ELBO penalises non-additivity, without these terms the model could inject
interaction structure into genuinely additive pairs at negligible likelihood cost, which would erase
precisely the additive signature that downstream GI quantification reads out. Finally, because
combinations reuse their parents' embeddings, gates, and the shared low-rank covariance $\Sigma$
rather than acquiring their own, a double perturbation that was never observed in training still
inherits a fully specified latent intervention---the property that makes out-of-distribution
combination prediction possible at all.

Under the structural-causal reading of the model (see \textbf{Theory}), a combination is an
intervention like any other: $\mathrm{do}(P{=}A{+}B)$ applied to a background abducted from
non-targeting control cells. Because our identifiability result is stated for a set of target
perturbations evaluated under a single abduction, the joint counterfactual over $\{A, B, A{+}B\}$---
and hence the interaction contrast $\Delta_{AB}-(\Delta_A+\Delta_B)$ defined on it---is identified
for any combination the assay contains. What is identified is the total shift $\Delta_{AB}$; its
split into parent scaling and interaction residual is not, and we therefore do not interpret the two
separately.

Identification is, however, relative to the distribution the model is fitted to, and this bounds what
can be claimed for combinations the model has not seen. Withholding a combination removes precisely
the conditional that identifies its mechanism, so the prediction the model then makes for it is not
an identified counterfactual but an extrapolation: it is produced by applying parent coefficients and
an interaction residual fitted on other pairs. What licenses it is structural rather than
inferential---because a combination reuses its constituents' embeddings, gates and covariance rather
than acquiring parameters of its own, it inherits a fully specified latent intervention even when
absent from training---and whether that extrapolation is accurate is an empirical question, not a
consequence of the identifiability result. We therefore do not claim to identify the counterfactual
of an unseen combination. We claim instead that the compositional parameterisation supports
extrapolation to one, and we test that claim directly: withheld combinations were assayed even though
they were not shown to the model, so their effects are independently measured and predictions for
them can be scored against measurement. The two tests are the held-out average treatment effect
below, and the held-out genetic-interaction classification in the following section.

Accordingly, we held out 20\% of the labelled double perturbations
stratified by GI class (22 of 88 pairs) and removed all of their cells from training, retaining every
single perturbation. We then quantified counterfactual average treatment effects (ATE) as the
Pearson correlation between predicted and observed log$_2$-CPM effect profiles, reported separately
for all perturbations, singles only, training-visible combinations, and held-out combinations
(Fig.~\ref{fig:fig6}b). Across five independently seeded runs SHERLOCK attained $r=0.852\pm0.004$
overall, $0.812\pm0.009$ on singles and $0.874\pm0.005$ on seen combinations, and---critically---
$0.836\pm0.005$ on combinations that were entirely absent from training, exceeding GEARS in every
column ($0.833\pm0.001$, $0.789\pm0.003$, $0.859\pm0.001$, $0.810\pm0.008$; Welch's $t$-test with
Benjamini--Hochberg correction across all comparisons, $q<0.003$ for every column). The generative
baselines are further behind in distribution (scGen $0.805\pm0.035$, cVAE $0.498\pm0.005$,
contrastiveVI $0.496\pm0.023$, sVAE $0.428\pm0.023$), scGen being the closest of them---SHERLOCK
exceeds it overall and on combinations ($q=0.042$ and $q=0.026$) though not on singles alone
($q=0.23$). More importantly, none of these baselines can be scored out of distribution at all: each
treats $A{+}B$ as an atomic label with its own independent embedding, so a combination absent from
training has no representation to evaluate and held-out ATE is undefined by construction rather than
merely poor. For fairness of comparison, the generative baselines were tuned with Optuna
\cite{akiba2019optuna}, SHERLOCK hyperparameters were selected by grid search, and GEARS was run with
the hyperparameters specified for the Norman dataset in its original publication
\cite{roohani2023gears}.

Because these conclusions rest on a learned latent geometry, we also quantified run-to-run stability
across the five seeds. The learned perturbation covariance was highly reproducible (leave-one-out
Spearman $r=0.788$, 95\% CI $[0.772,0.804]$ between each run's $\rho$ correlation matrix and the mean
of the others), as were the counterfactual effect magnitudes ($r=0.924\pm0.007$ between predicted and
observed per-perturbation effect norms), while the raw gate co-activation pattern was more variable
($r=0.494$, 95\% CI $[0.446,0.533]$)---expected, since individual latent coordinates are exchangeable
across runs and only their induced relational structure is identified
(Supplementary Fig.~\ref{fig:SuppNormanRepro}).

The structure that makes this composition work is visible directly in the learned perturbation
covariance $\Sigma$, which is shared between singles and combinations and is therefore the object
through which a combination inherits its parents' mechanisms. Its correlation structure recovers
gene families and signalling relationships without any pathway supervision
(Supplementary Fig.~\ref{fig:SuppNormanCov}). The most strongly coupled perturbation pairs are
almost uniformly paralogues or direct pathway partners: the posterior \textit{HOX13} paralogues
\textit{HOXA13}/\textit{HOXC13} ($r=0.89$), \textit{ETS2}/\textit{MAPK1} ($r=0.89$), where ETS2 is a
canonical ERK substrate, the CIP/KIP cyclin-dependent kinase inhibitors \textit{CDKN1A},
\textit{CDKN1B} and \textit{CDKN1C} (pairwise $r=0.79$--$0.89$) \cite{sherr1999cdk}, the tyrosine
phosphatases \textit{PTPN1}, \textit{PTPN9} and the TULA-family phosphatase \textit{UBASH3B}
($r=0.84$--$0.88$) \cite{tsygankov2013tula}, the C/EBP bZIP factors \textit{CEBPA}/\textit{CEBPB}
($r=0.87$) \cite{friedman2007cebp}, the forkhead paralogues \textit{FOXA1}/\textit{FOXA3} ($r=0.85$)
and \textit{FOXF1}/\textit{FOXL2} ($r=0.76$), the two p38 MAP2Ks \textit{MAP2K3}/\textit{MAP2K6}
($r=0.82$) \cite{cuenda2007p38}, and the immunoglobulin-domain paralogues
\textit{IGDCC3}/\textit{PRTG} ($r=0.81$). Quantitatively, every perturbation module annotated by
Norman et al. (Supplementary Table 6 of \cite{norman2019exploring}) is enriched in the learned
covariance relative to the rest of the screen: mean within-module versus between-module correlation
is $0.82$ versus $-0.20$ for the MAP2K module, $0.81$ versus $0.12$ for the Ig-domain module, $0.66$
versus $-0.01$ for the bZIP transcription-factor module, $0.65$ versus $0.12$ for the tyrosine
phosphatases, $0.63$ versus $0.21$ for the MAP4K module, and $0.50$ versus $0.20$ for the
regulation-of-apoptosis module. Hierarchical clustering of $\Sigma$ additionally isolates coherent
programs that cut across those annotations---an erythroid/Notch block
(\textit{KLF1}, \textit{CSRNP1}, \textit{MAML2}, \textit{TGFBR2}, \textit{ZC3HAV1}), with KLF1 the
master erythroid regulator \cite{siatecka2011klf1}; a myeloid-differentiation block containing the
three C/EBP factors together with \textit{SPI1} (PU.1) and the ETO-family corepressors
\textit{RUNX1T1} and \textit{CBFA2T3} \cite{rosenbauer2007transcription,davis2003eto}; a p38/MAPK
block (\textit{MAP2K3}, \textit{MAP2K6}, \textit{ELMSAN1}, \textit{PTPN12}, \textit{SET},
\textit{TBX2}, \textit{IRF1}); and a two-member apoptotic block pairing the pore-forming effector
\textit{BAK1} with the BH3-only activator \textit{BCL2L11} (BIM) \cite{youle2008bcl2}. The remaining
majority of perturbations forms a large weakly coupled group, consistent with the observation that
most single CRISPRa perturbations in this screen produce only modest transcriptional responses
\cite{norman2019exploring}.

These covariance modules are what the SHERLOCK latent space renders geometrically
(Fig.~\ref{fig:fig6}a). Restricting to the Norman Supplementary Table 6 modules and to every
combination containing at least one module member (45,231 cells, 116 perturbations), the annotated
modules---MAP2K, MAP4K, tyrosine phosphatase, Ig-domain-containing, kinesin, bZIP transcription
factor, homeobox transcription factor, Brachyury transcription factor, and regulation of
apoptosis---occupy coherent territories, and combinations of two module members are placed in the
region defined by their parents rather than in isolated islands
(Supplementary Fig.~\ref{fig:SuppNormanLatent}). Perturbations that share a covariance module also
share latent territory, so the colouring of the embedding and the block structure of $\Sigma$ are two
views of the same learned quantity: mechanistic relatedness of the parents is inherited by the pair.
This is the behaviour expected of a modular genetic-interaction landscape, in which functionally
related regulators produce convergent transcriptional outcomes and interactions concentrate within
and between coherent functional modules \cite{norman2019exploring,horlbeck2018mapping}.

Individual pairs make the correspondence between latent geometry and interaction type explicit. For
\textit{IKZF3}$+$\textit{MAP2K6}, annotated as epistatic, the double-perturbed cells overlay the
\textit{MAP2K6} single population and are clearly displaced from the \textit{IKZF3} population---the
direct latent signature of one parent's program masking the other. This is the expected direction
biologically: MAP2K6 is a dedicated activator of p38 MAPK and drives a strong stress and
differentiation response in K562 \cite{cuenda2007p38}, whereas IKZF3 (Aiolos) is a lymphoid Ikaros
family factor with limited transcriptional consequence in this myeloid background. The corresponding
GI metrics are strongly asymmetric (dominance $0.54$, parent distance-correlation difference $0.52$,
equality of contribution $0.44$) and SHERLOCK assigns the pair to epistasis, matching the reference
label. For \textit{CEBPE}$+$\textit{SET}, the double-perturbed cells sit markedly closer to the
unperturbed background than either single, with a joint effect less than half the additive
expectation (ratio of joint to additive norm $0.44$) despite a well-fit parent regression---the
signature of mutual suppression, and the class SHERLOCK assigns
(Supplementary Fig.~\ref{fig:SuppNormanLatent}). This too is consistent with the antagonism of the
two programs involved: CEBPE drives granulocytic differentiation \cite{friedman2007cebp}, whereas SET
(I2PP2A) is an oncoprotein inhibitor of PP2A associated with a differentiation block in myeloid
leukaemia, and every other C/EBP pairing that carries a reference label in this screen falls into the
redundancy or suppression family.

\subsection{SHERLOCK recovers Norman genetic-interaction categories from decoded perturbation effects}

We next asked whether the modelled combination effects are quantitative enough to reproduce the
expert GI assignments of Norman et al. We adopted their own formulation: for each double
perturbation, the joint effect profile is regressed on its two single-perturbation profiles by
Theil--Sen regression without intercept \cite{sen1968estimates}, yielding coefficients $c_1,c_2$ from
which magnitude $\sqrt{c_1^2+c_2^2}$, dominance $|\log_{10}(|c_1|/|c_2|)|$, model fit, equality of
contribution, and the distance correlation between the two singles are derived, together with
residual-geometry terms comparing the joint effect to the additive expectation (the additive residual
norm, the ratio of joint to additive norm, and the signed alignment of the additive residual with the
joint effect). These quantities map onto the biological definitions directly: a dominant coefficient
with unequal parental contribution is epistasis; near-identical single profiles whose combination
adds little is redundancy; a joint response smaller or larger than the additive expectation while
retaining parental direction is suppression or potentiation; and a joint response poorly explained by
any weighting of its parents is neomorphic. Categories are assigned by an equal-weight composite of
$z$-scored metrics---one score per class, argmax over classes. Crucially, this entire downstream is a
single shared routine applied identically to SHERLOCK and to GEARS; the only difference between the
two evaluations is which per-perturbation effect vectors enter it, so the comparison isolates the
quality of the modelled effects rather than the scoring convention. The Norman annotation splits
synergy into similar- and dissimilar-phenotype subtypes; we score synergy as a single class, giving
seven categories evaluated over the 86 labelled pairs for which both models produce complete metrics.

The main-figure evaluation takes both the single- and the double-perturbation effect vectors from
the model's decode of the posterior $q(z)$ over each perturbation's own cells, referenced to a common
decoded control baseline. Two considerations motivate this over the composed-shift contrast, which by
Corollary~\ref{cor:gi-id} targets the same estimand. First, generating all three vectors from the
composed shift makes the regression of the double on its singles close to degenerate by construction:
the fitted coefficients and residual would return the model's own parent coefficients and interaction
term rather than a comparison against data. Second, the Norman categories were assigned from measured
joint profiles, so an estimator whose joint effect derives from the cells that received the pair
mirrors the construction the labels came from.
In this setting SHERLOCK outperforms GEARS on every summary metric (Fig.~\ref{fig:fig6}d): weighted
F1 $0.712\pm0.017$ versus $0.623\pm0.035$, adjusted Rand index $0.440\pm0.017$ versus $0.387\pm0.033$,
and adjusted mutual information $0.512\pm0.022$ versus $0.446\pm0.050$ across five runs (Welch's
$t$-test with Benjamini--Hochberg correction across the three metrics; $q=0.008$, $0.028$ and $0.040$
respectively). Per-class one-versus-rest AUC-ROC shows that the advantage is not carried by a single
easy category: SHERLOCK reaches $0.979\pm0.005$ for epistasis, $0.860\pm0.038$ for suppression,
$0.845\pm0.009$ for synergy, $0.810\pm0.005$ for potentiation, $0.770\pm0.021$ for approximately
additive, $0.767\pm0.028$ for neomorphic and $0.736\pm0.027$ for redundancy, exceeding GEARS on five
of the seven classes. The largest margins fall on the two categories that depend on the additive null
being represented faithfully---approximately additive ($0.770$ versus $0.657$) and suppression
($0.860$ versus $0.713$)---which is the behaviour the additivity prior on $c$ and $s$ was introduced
to protect. GEARS is modestly stronger on potentiation and redundancy. Embedding the pairs in the
space of GI metrics gives the same picture qualitatively: the SHERLOCK-derived embedding separates
the Norman classes more cleanly than the GEARS-derived one (macro-averaged silhouette in the metric
space $0.085$ versus $0.050$; Fig.~\ref{fig:fig6}c and Supplementary
Fig.~\ref{fig:SuppNormanGears}), with suppression, epistasis and approximately additive forming the
most coherent territories and neomorphic---by definition a residual category of responses unlike
either parent---the most diffuse.

Two controls delimit what this result does and does not claim. First, we repeated the evaluation on
the 22 held-out pairs with each pair's effect generated from the composed shift, so that both models
predict a combination from its two constituents and neither uses the pair's own cells. This is the
second of the two tests of compositional extrapolation set out above, now at the level of interaction
type rather than of effect magnitude: the categories are assigned from a joint effect the model
constructed rather than measured, and scored against the annotation derived from the measured one.
The two models are indistinguishable on it (SHERLOCK weighted F1 $0.692\pm0.047$, GEARS
$0.696\pm0.120$; $q=0.95$ for all three summary metrics after correction; Supplementary
Fig.~\ref{fig:SuppNormanGIheldout}). With 22 pairs distributed over seven classes the per-class
supports are in the low single digits, and the GEARS run-to-run spread is more than twice
SHERLOCK's, so this test resolves only large differences; we report it as evidence that composition
extrapolates to unseen combinations as well as a graph-based predictor does, not as evidence of
superiority out of distribution. Second, we repeated the GI
evaluation protocol of the GEARS study---their published metric and threshold definitions, scored by
precision@10 and F1 over the five original Norman subtypes---using observed (raw) expression profiles
as an additional reference alongside both models (Supplementary Fig.~\ref{fig:SuppNormanGIraw}). Raw
expression is the strongest source in that protocol (overall F1 $0.573$; precision@10 of $0.9$ for
neomorphism and $0.7$ for epistasis, against $0.34$ and $0.28$ for the best model-derived scores).
This is expected and, we would argue, appropriate: the Norman categories are defined \emph{on}
measured expression, so a procedure that consumes measured expression directly is scoring against its
own definition, including whatever unmodelled noise structure that expression carries. We do not
claim to improve on genetic interactions computed from raw expression, and we do not reformulate them
in latent space---which the model could readily do, but which would define the categories in a space
different from the one in which the reference labels were assigned. A representation-learning model
is not well served by being judged on absolute expression reconstruction; the relevant question is
whether a compact, interpretable latent parameterisation retains enough interaction structure to
recover the expert categories, and the answer here is that it does.

Finally, we asked whether the assigned categories correspond to interpretable transcriptional
patterns rather than only to positions in a derived metric space. For each Norman category we
selected a correctly classified held-out pair and reconstructed the effect profiles of the two
singles and the double from the model's own decode, restricted to the 100 highest-variance genes
across perturbations and hierarchically ordered (Fig.~\ref{fig:fig6}e). All seven exemplars are drawn
from the held-out set, so each displayed reconstruction is a genuine extrapolation, and each matches
its category both visually and mechanistically. The redundant exemplar
\textit{MAP2K3}$+$\textit{MAP2K6} pairs the two MAP2Ks that phosphorylate p38 on the same activation
loop residues \cite{cuenda2007p38}; their single profiles are nearly identical (singles similarity
$0.94$) and the combination adds little beyond either ($c$-magnitude $1.05$, joint-to-additive norm
ratio $0.80$)---textbook functional redundancy, and the same relationship the covariance matrix
independently assigns the highest within-module correlation in the screen. The suppression exemplar
\textit{CEBPB}$+$\textit{PTPN12} combines a myeloid differentiation driver with a tyrosine phosphatase
that dampens receptor tyrosine kinase signalling, and yields a joint response attenuated relative to
both parents. The synergy exemplar \textit{CBL}$+$\textit{CNN1} shows a joint response exceeding the
additive expectation while retaining the direction of both parents (joint-to-additive ratio $1.23$,
singles similarity $0.84$), consistent with CBL acting as an E3 ligase that terminates
growth-factor-receptor signalling \cite{mohapatra2013cbl} and with the erythroid-differentiation
program Norman et al. report for CBL-containing combinations; the closely related pair
\textit{CBL}$+$\textit{PTPN12}, which joins two negative regulators of the same signalling axis, is
likewise correctly called as synergy. The epistasis exemplar
\textit{CEBPA}$+$\textit{ZC3HAV1} is dominated by the C/EBP program, and the neomorphic exemplars
\textit{KLF1}$+$\textit{TGFBR2} and \textit{DUSP9}$+$\textit{MAPK1} produce responses poorly explained
by any weighting of their parents---in the latter case unsurprisingly, since DUSP9 is a dual-specificity
phosphatase that directly inactivates ERK2/MAPK1, so co-activating a kinase and its own phosphatase
places the cell in a state neither single perturbation visits \cite{caunt2013dual}.

The remaining 45 double perturbations in the screen carry no reference annotation in the label set
available to us, and the model assigns each of them a category as well. We treat these as internal
consistency checks rather than as findings: several fall where the pair's own covariance module would
place it, including \textit{CDKN1A}$+$\textit{CDKN1B} (p21/p27) called redundant, as expected for two
CIP/KIP inhibitors acting on overlapping cyclin--CDK complexes \cite{sherr1999cdk};
\textit{UBASH3A}$+$\textit{UBASH3B} called synergistic, matching the paralogous
TULA-family phosphatases \cite{tsygankov2013tula}; and \textit{BAK1}$+$\textit{BCL2L11}, the pair that
forms its own two-member apoptotic block in the learned covariance \cite{youle2008bcl2}. That the
unlabelled calls agree with the relationships the covariance matrix encodes, without any of these
pairs contributing to the labelled evaluation, supports the interpretation that the compositional
latent parameterisation captures interaction structure rather than fitting the labelled subset.

Taken together, these results show that representing a combination as a structured composition of its
constituents' latent mechanisms---rather than as an independent embedding---is what makes both
out-of-distribution combination prediction and quantitative genetic-interaction recovery possible
within a single interpretable generative model.